\documentclass[11pt, letterpaper]{article}
\usepackage[utf8]{inputenc}
\usepackage[T1]{fontenc}
\usepackage{amsmath}
\usepackage{amsfonts}
\usepackage{amssymb}
\usepackage{geometry}
\usepackage{verbatim}

\geometry{letterpaper, margin=1in}

\title{\textbf{Problem C: Jayden's T-Spin Setup}}
\author{Time Limit: 1.0s | Memory Limit: 256MB}
\date{}

\begin{document}

\maketitle

\section*{Description}

Jayden is an S-rank player on \textit{Tetrio}. He is currently practicing his opening setups. To perfect his T-Spin technique, he needs to simulate how pieces stack up on the board.

The game board is a grid with a width of \textbf{10 columns} and a maximum height of \textbf{20 rows}. The rows are numbered 1 to 20 from bottom to top. The columns are numbered 1 to 10 from left to right.

Jayden drops $N$ pieces into the grid one by one. For this practice session, the pieces do \textbf{not} rotate. They always appear in their default "spawn" orientation (described below) and fall straight down from the top of the specified column until they land on the floor (Row 0) or another piece.

The pieces are the standard Tetrominoes, defined as follows (using `\#' for blocks and `.' for empty space):

\begin{verbatim}
Type I:    Type O:    Type T:    Type L:    Type J:    Type S:    Type Z:
####       ##         .#.        ..#        #..        .##        ##.
           ##         ###        ###        ###        ##.        .##
\end{verbatim}

\textbf{Note on Input:}
You will be given the type of the piece and the column index of its \textbf{left-most} block. The piece falls until any part of it collides with a block below it or the floor.

If any part of a piece settles at a height greater than 20, the board overflows, and the game ends immediately.

Your task is to process the sequence of moves. If the game ends due to overflow, output ``Game Over''. Otherwise, output the height of the highest occupied block on the board after all pieces have been placed.

\section*{Input}

The first line contains a single integer $N$ ($1 \le N \le 1000$), the number of pieces.

The next $N$ lines each describe a move with a character $T$ and an integer $C$:
\begin{itemize}
    \item $T \in \{'I', 'O', 'T', 'L', 'J', 'S', 'Z'\}$ is the piece type.
    \item $C$ ($1 \le C$) is the column index of the left-most column of the piece.
\end{itemize}
It is guaranteed that the piece fits horizontally (e.g., an 'I' piece of width 4 will not be dropped at column 8, 9, or 10).

\section*{Output}

If the stack height exceeds 20 at any point, output ``Game Over''.
Otherwise, output a single integer representing the maximum row index occupied.

\section*{Examples}

\subsection*{Example 1}
\textbf{Input:}
\begin{verbatim}
3
I 1
I 1
O 5
\end{verbatim}

\textbf{Output:}
\begin{verbatim}
2
\end{verbatim}

\textbf{Explanation:} \\
1. An 'I' piece (flat, height 1) is dropped at column 1. It lands on row 1. (Occupies rows 1, cols 1-4).
2. Another 'I' piece is dropped at column 1. It lands on top of the first. (Occupies row 2, cols 1-4).
3. An 'O' piece (2x2) is dropped at column 5. It lands on row 1-2. (Occupies rows 1-2, cols 5-6).
Max height is 2.

\subsection*{Example 2}
\textbf{Input:}
\begin{verbatim}
4
T 1
Z 2
J 1
I 1
\end{verbatim}

\textbf{Output:}
\begin{verbatim}
5
\end{verbatim}

\end{document}
